\section{この付録で見るもの}

本編では,漸化式を与えられたときに,その一般項を求めるための考え方を整理してきた.
第1章から始まり,第2章では基本となる知識を,第3章では,漸化式を分類しそれぞれに対してのアプローチを扱った.
第4章ではそれらをまとめ上げ,本書の主題である"体系的解説"を行った.

付録では,問題を解くためのパターンを追加するのではなく,これまでの解法がより高度な数学で,あるいは別の視点から見て,どのような理論を内包しているかを眺める.
一つの漸化式を少し遠くから見ると,次のような問いが現れる.

\begin{itemize}
  \item 隣り合う項の差を,微分や積分に似た操作として扱えないか.
  \item $a_n$を無限に遠くまで見通すと最終的にどうなるのか.
  \item $a_n$の$n$を実数に拡張するとどうなるのか.
  \item 度々登場した特性方程式と不動点をより統一的に扱えないか.
\end{itemize}

これらの疑問をより広い数学の世界から見渡して,そこから生まれる新たな数学の道具が生まれる流れを追うものである.
その道具はしばしば大学数学を前提としたものになる.
しかし,付録を読むために,すべての大学数学を先に知っている必要はない.
あくまで順番に知的好奇心を満たすために,疑問を解決していくだけである.

各節では,まず問題意識を示し,次に必要な用語を定義し,最後に漸化式との
つながりを確認する.証明を完全に一般化することよりも,なぜその概念が
必要になり,何を見えるようにしたのかを大切にする.